\documentclass[11pt]{article}
\usepackage[margin=0.7in]{geometry}
\usepackage{amsmath,amssymb}
\usepackage{enumitem}
\usepackage{tikz}
\usepackage{pgfplots}
\pgfplotsset{compat=1.18}
\usepackage{fancyhdr}
\usepackage{parskip}
\setlist[enumerate]{leftmargin=*, itemsep=6pt}
\pagestyle{fancy}
\fancyhf{}
\fancyfoot[L]{\small\textit{en}Vision\textsuperscript{\textregistered} Algebra 2}
\fancyfoot[C]{\small 2-6 Lesson Quiz, short form}
\fancyfoot[R]{\small\thepage}
\renewcommand{\headrulewidth}{0pt}
\newcommand{\choice}[1]{\tikz[baseline=-0.6ex]\node[draw, circle, inner sep=0.6pt, minimum size=3.3mm, font=\scriptsize\bfseries] {#1};\;}
\newcommand{\checkbox}{$\square$\;}

\begin{document}

\noindent\textbf{Name}\; \rule{6cm}{0.4pt} \hfill Period \rule{1.2cm}{0.4pt} \quad Date \rule{2.2cm}{0.4pt}

\vspace{4pt}\noindent\rule{\textwidth}{0.6pt}

\vspace{2pt}\noindent{\LARGE\textbf{2-6 \; Lesson Quiz: The Quadratic Formula}} \hfill {\small short form}

\vspace{4pt}\noindent\rule{\textwidth}{0.6pt}

\vspace{6pt}
\begin{minipage}[t]{0.50\textwidth}
\textbf{1.} Solve $x^2+8x-6=0$ using the Quadratic Formula.

\vspace{4pt}
\begin{itemize}[leftmargin=1.2em, itemsep=5pt]
\item[] \choice{A} $x=32+2\sqrt{2}$ and $x=32-2\sqrt{2}$
\item[] \choice{B} $x=-4+\sqrt{22}$ and $x=-4-\sqrt{22}$
\item[] \choice{C} $x=12$ and $x=-\frac{1}{2}$
\item[] \choice{D} $x=4+2\sqrt{22}$ and $x=4-2\sqrt{22}$
\end{itemize}
\end{minipage}\hfill
\begin{minipage}[t]{0.46\textwidth}
Show your work: $a=\rule{0.8cm}{0.4pt}$, $b=\rule{0.8cm}{0.4pt}$, $c=\rule{0.8cm}{0.4pt}$

\vspace{14pt}\rule{\linewidth}{0.4pt}

\vspace{14pt}\rule{\linewidth}{0.4pt}

\vspace{14pt}\rule{\linewidth}{0.4pt}

\vspace{14pt}\rule{\linewidth}{0.4pt}
\end{minipage}

\vspace{10pt}\noindent\rule{\textwidth}{0.4pt}

\vspace{6pt}
\begin{minipage}[t]{0.50\textwidth}
\textbf{2.} Solve $x^2+2x+7=0$ using the Quadratic Formula. Select \textbf{all} solutions that apply.

\vspace{4pt}
\begin{itemize}[leftmargin=1.2em, itemsep=5pt]
\item[] \checkbox \textbf{A.} $x=-1+i\sqrt{6}$
\item[] \checkbox \textbf{B.} $x=-1+\sqrt{6}$
\item[] \checkbox \textbf{C.} $x=-1-i\sqrt{6}$
\item[] \checkbox \textbf{D.} $x=-1-\sqrt{6}$
\item[] \checkbox \textbf{E.} $x=1-\sqrt{6}$
\end{itemize}
\end{minipage}\hfill
\begin{minipage}[t]{0.46\textwidth}
Show your work:

\vspace{14pt}\rule{\linewidth}{0.4pt}

\vspace{14pt}\rule{\linewidth}{0.4pt}

\vspace{14pt}\rule{\linewidth}{0.4pt}

\vspace{14pt}\rule{\linewidth}{0.4pt}
\end{minipage}

\vspace{10pt}\noindent\rule{\textwidth}{0.4pt}

\vspace{6pt}
\begin{minipage}[t]{0.50\textwidth}
\textbf{3.} Describe the nature of the roots for the equation $49x^2-28x+4=0$.

\vspace{4pt}
\begin{itemize}[leftmargin=1.2em, itemsep=5pt]
\item[] \choice{A} two real roots
\item[] \choice{B} one real root
\item[] \choice{C} two complex roots
\item[] \choice{D} one complex root
\end{itemize}
\end{minipage}\hfill
\begin{minipage}[t]{0.46\textwidth}
Show the discriminant: \ $b^2-4ac=$

\vspace{14pt}\rule{\linewidth}{0.4pt}

\vspace{14pt}\rule{\linewidth}{0.4pt}

\vspace{8pt}
Explain your choice in one sentence using the word \emph{discriminant}:

\vspace{14pt}\rule{\linewidth}{0.4pt}
\end{minipage}

\vfill
\noindent\textbf{Scoring:} Question 1 and Question 3 are 10 points each. Question 2 is a bonus of up to 5 points (both correct boxes, none wrong). Total out of 20.

\newpage
\noindent{\LARGE\textbf{2-6 Lesson Quiz, short form: answer key}}

\vspace{4pt}\noindent\rule{\textwidth}{0.6pt}

\begin{enumerate}
\item \textbf{B, $x=-4\pm\sqrt{22}$.} $a=1$, $b=8$, $c=-6$: $x=\dfrac{-8\pm\sqrt{64+24}}{2}=\dfrac{-8\pm\sqrt{88}}{2}=\dfrac{-8\pm2\sqrt{22}}{2}=-4\pm\sqrt{22}$. Savvas quiz item 1, DOK 2, Realize skills A31, A72 (Example 1A). \textbf{10 points:} 6 for B, 4 for $a$, $b$, $c$ and a correct substitution.
\item \textbf{A and C, $x=-1\pm i\sqrt{6}$.} $a=1$, $b=2$, $c=7$: $x=\dfrac{-2\pm\sqrt{4-28}}{2}=\dfrac{-2\pm\sqrt{-24}}{2}=\dfrac{-2\pm2i\sqrt6}{2}=-1\pm i\sqrt6$. B and D drop the $i$ (the radicand $-24$ is negative); E has the wrong sign on $-1$ and also drops the $i$. Savvas quiz item 2, DOK 2, Realize skills A31, A72 (Example 1B). \textbf{Bonus 5:} A and C and no other box; 3 for one of A or C and no wrong box; 0 otherwise.
\item \textbf{B, one real root.} $b^2-4ac=(-28)^2-4(49)(4)=784-784=0$, so there is one real root ($x=\tfrac{28}{98}=\tfrac27$; the equation is $(7x-2)^2=0$). Savvas quiz item 3, DOK 1, Realize skill A32 (Example 3). \textbf{10 points:} 6 for B, 2 for the discriminant $0$, 2 for the sentence naming the discriminant.
\end{enumerate}

\textbf{Items left off this form} (available for a retake or Form B):
Savvas item 4, Richard's ball, $h(t)=-5t^2+10t+6$: does it reach 12 m? (answer D, the ball does not reach 12 m: $-5t^2+10t-6=0$ has discriminant $100-120=-20<0$; DOK 2, A31, A72; taught by Example 4, cut, and rehearsed by the packet's Practice page);
Savvas item 5, values of $b$ that give $27x^2+bx+3=0$ one real solution (answer A and F, $b=\pm18$: $b^2-4(27)(3)=b^2-324=0$; DOK 2, A32; taught by Example 5, cut).

\textbf{IXL replacement:} SmartScore 70 or above on YQH and QHK replaces this quiz with 20/20.

\end{document}
